\documentclass[11pt]{article}
\usepackage[a4paper,margin=1.9cm]{geometry}
\usepackage[T1]{fontenc}
\usepackage{textcomp}
\usepackage[spanish]{babel}
\usepackage{amsmath,amssymb}
\usepackage{lmodern}
\renewcommand{\familydefault}{\sfdefault}
\usepackage{enumitem}
\newcommand{\vect}[2]{\begin{pmatrix}#1\\#2\end{pmatrix}}
\usepackage[table]{xcolor}
\usepackage{parskip}
\definecolor{unalblue}{RGB}{31,73,124}
\newlist{opciones}{enumerate}{1}
\setlist[opciones]{label=\alph*),itemsep=4pt,topsep=4pt,leftmargin=1.8em,font=\normalfont}
\setlist[enumerate,1]{itemsep=1em,topsep=0.8em}
\newcommand{\examfooter}{\vfill\rule{\linewidth}{0.4pt}\\[1pt]{\scriptsize\textit{Transcripción digital --- MathUNAL}\hfill\textcolor{unalblue}{\textbf{\thepage}}}}
\pagestyle{empty}

\begin{document}
\noindent
\begin{minipage}[t]{0.6\linewidth}
{\small\bfseries Geometría Vectorial y Analítica --- Parcial 3}\\
{\small Semestre 2022--01}
\end{minipage}%
\hfill
\begin{minipage}[t]{0.35\linewidth}
\raggedleft\small Escuela de Matemáticas. Universidad Nacional de Colombia, Sede Medellín.
\end{minipage}\\[2pt]
\rule{\linewidth}{0.6pt}

\begin{center}
\textbf{Geometría Vectorial --- Parcial 3}\\
Junio del 2022
\end{center}

\textbf{Instrucciones:} La duración del examen es \textbf{1:50 horas}. Verifique que sus datos de identificación son correctos. La prueba consta de cinco preguntas. Escriba sus respuestas a los ejercicios 1,2,3,4 en la tabla de puntaje. No use fracciones, solamente decimales y redondee a dos cifras decimales. La pregunta 5 es de procedimiento, y su respuesta debe escribirse debajo del enunciado. No está permitido: hacer preguntas, usar calculadoras ni trabajar en hojas diferentes a las del examen. No se permite sacar el teléfono celular durante el examen.

\begin{center}
\begin{tabular}{|c|c|c|}
\hline
Ejercicio & Porcentaje & Respuesta\\
\hline
1 & 20\% & \\
\hline
2 & 16\% & \\
\hline
3 & 16\% & \\
\hline
4 & 18\% & \\
\hline
5 & 30\% & Procedimiento\\
\hline
\end{tabular}
\end{center}

\begin{enumerate}
\item Sean $A=\vect{2}{7/6}$, $B=\vect{1+1/6}{-2}$, y $C=\vect{-1}{-7/12}$. De las siguientes afirmaciones, seleccione todas aquellas que son verdaderas.
\begin{opciones}
\item Cualquier otro vector de $\mathbb{R}^2$ se puede escribir como una combinación lineal de $B$ y $C$.
\item El vector $C-B$ está a la derecha de $A$.
\item $A$ y $C$ son linealmente independientes.
\item $A$ está en la recta que pasa por $B$ y $C$.
\item $C^\perp$ pertenece a la recta generada por $B$: $\mathcal{L}_B$.
\item Ninguna de las anteriores.
\end{opciones}

\item Supongamos que la proyección ortogonal del vector $X=\vect{11}{2}$ sobre el vector $U\neq 0$ es el vector $\vect{3}{-4}$. Entonces, la magnitud de la proyección ortogonal de $X$ sobre el vector $U^\perp$ es:

\item Considere los vectores $A=\vect{1}{2}$, $B=\vect{1}{-1}$, y $C=\vect{1}{1}$. Encuentre el vector $D$ tal que $[\overrightarrow{AB}]+[\overrightarrow{AC}]=[\overrightarrow{AD}]$.

\item Considere los puntos $A=\vect{-2}{3}$, $B=\vect{0}{5}$, $C=\vect{-5}{7}$. De las siguientes afirmaciones, seleccione todas aquellas que son verdaderas.
\begin{opciones}
\item $\tan\angle(BAC)=7$.
\item $\cos\angle(BAC)=\dfrac{1}{5\sqrt{2}}$.
\item $\sin\angle(CAB)=\dfrac{7\sqrt{2}}{10}$.
\item $\sin\angle(BAC)=\dfrac{7}{5\sqrt{2}}$.
\item $\tan\angle(CAB)=1/7$.
\item $\cos\angle(CAB)=\dfrac{1}{5\sqrt{2}}$.
\item Ninguna de las anteriores.
\end{opciones}

\item \textbf{Pregunta de procedimiento:} En el espacio de abajo escriba un procedimiento completo justificando cada paso.

Considere los puntos $A=\vect{5}{3}$ y $B=\vect{q}{1}$. Halle $q$ para que el ángulo desde $A$ hasta $B$, medido en sentido antihorario, sea de $135^\circ$.
\end{enumerate}

\examfooter
\end{document}
